\frfilename{mt257.tex}
\versiondate{14.8.13}
\copyrightdate{1995}

\def\chaptername{Ukuran produk}
\def\sectionname{Konvolusi ukuran}

\newsection{257}

Gagasan-gagasan dalam bab ini dapat dipadukan secara memuaskan dalam
teori konvolusi ukuran Radon,
yang akan berguna dalam \S272 dan sekali lagi dalam
\S285.   Di sini saya hanya memberikan definisi (257A) dan sifat utama
(257B) dari konvolusi ukuran Radon berhingga total, beserta beberapa
akibat dan catatan mengenai hubungan antara konvolusi fungsi
dan konvolusi ukuran (257F).

\leader{257A}{Definisi}  Misalkan $r\ge 1$ suatu
bilangan bulat dan $\nu_1$, $\nu_2$ dua ukuran Radon berhingga total pada
$\Bbb R^r$.   Misalkan $\lambda$ adalah ukuran produk pada
$\BbbR^r\times\BbbR^r$\cmmnt{;   maka $\lambda$ juga merupakan ukuran Radon
(berhingga total), berdasarkan 256K}.   Definisikan
$\phi:\BbbR^r\times\BbbR^r\to\BbbR^r$ dengan menetapkan
$\phi(x,y)=x+y$\cmmnt{;  maka
$\phi$ kontinu, sehingga terukur dalam pengertian
256G}.   {\bf Konvolusi}
dari $\nu_1$ dan $\nu_2$, yaitu $\nu_1*\nu_2$, adalah ukuran citra
$\lambda\phi^{-1}$;  \cmmnt{berdasarkan 256G,} ini merupakan ukuran Radon.

Perhatikan bahwa jika $\nu_1$ dan $\nu_2$ adalah ukuran probabilitas Radon, maka
$\lambda$ dan $\nu_1*\nu_2$ juga merupakan ukuran probabilitas.

\leader{257B}{Teorema} Misalkan $r\ge 1$ suatu bilangan bulat, dan $\nu_1$ serta
$\nu_2$ dua ukuran Radon berhingga total pada $\BbbR^r$;  misalkan
$\nu=\nu_1*\nu_2$ adalah konvolusi keduanya, dan $\lambda$ produk keduanya pada
$\BbbR^r\times\BbbR^r$.   Maka untuk setiap fungsi bernilai real $h$ yang didefinisikan
pada suatu himpunan bagian dari $\BbbR^r$,

\Centerline{$\int h(x+y)\lambda(d(x,y))$ ada dan nilainya $=\int h(x)\nu(dx)$}

\noindent jika salah satu integral terdefinisi dalam $[-\infty,\infty]$.
%will this work for $[-\infty,\infty]$-valued h?

\proof{ Terapkan 235J dengan $J(x,y)=1$, $\phi(x,y)=x+y$ untuk semua $x$,
$y\in\BbbR^r$.
}%end of proof of 257B

\leader{257C}{Akibat} Misalkan $r\ge 1$ suatu bilangan bulat, dan $\nu_1$,
$\nu_2$ dua ukuran Radon berhingga total pada $\BbbR^r$;  misalkan
$\nu=\nu_1*\nu_2$ adalah konvolusi keduanya, dan $\lambda$ produk keduanya pada
$\BbbR^r\times\BbbR^r$;  tuliskan $\Sigma$ untuk domain $\nu$.
Misalkan $h$ adalah fungsi yang terukur terhadap $\Sigma$
dan terdefinisi $\nu$-hampir di mana-mana pada $\BbbR^r$.     Andaikan salah
satu dari integral-integral

\Centerline{$\iint|h(x+y)|\nu_1(dx)\nu_2(dy)$,
\quad$\iint|h(x+y)|\nu_2(dy)\nu_1(dx)$,
\quad$\int h(x+y)\lambda(d(x,y))$}

\noindent ada dan berhingga.   Maka $h$ terintegralkan terhadap $\nu$ dan

\Centerline{$\int h(x)\nu(dx)=\iint h(x+y)\nu_1(dx)\nu_2(dy)
=\iint h(x+y)\nu_2(dy)\nu_1(dx)$.}

\proof{ Gabungkan 257B dengan teorema Fubini dan Tonelli (252H).
}%end of proof of 257C

\vleader{48pt}{257D}{Akibat} Jika $\nu_1$ dan $\nu_2$ adalah ukuran Radon berhingga total
pada $\BbbR^r$, maka $\nu_1*\nu_2=\nu_2*\nu_1$.

\proof{ Untuk setiap himpunan Borel $E\subseteq\BbbR^r$, terapkan 257C pada $h=\chi E$
untuk memperoleh

$$\eqalign{(\nu_1*\nu_2)(E)
&=\iint\chi E(x+y)\nu_1(dx)\nu_2(dy)
=\iint\chi E(x+y)\nu_2(dy)\nu_1(dx)\cr
&=\iint\chi E(y+x)\nu_2(dy)\nu_1(dx)
=(\nu_2*\nu_1)(E).\cr}$$

\noindent Jadi $\nu_1*\nu_2$ dan $\nu_2*\nu_1$ bernilai sama pada himpunan-himpunan Borel
di $\BbbR^r$;  karena keduanya merupakan ukuran Radon, keduanya harus
identik (256D).
}%end of proof of 257D

\leader{257E}{Akibat} Jika $\nu_1$, $\nu_2$ dan $\nu_3$ adalah ukuran Radon
berhingga total pada $\BbbR^r$, maka
$(\nu_1*\nu_2)*\nu_3=\nu_1*(\nu_2*\nu_3)$.

\proof{ Untuk setiap himpunan Borel $E\subseteq\BbbR^r$, terapkan 257B pada $h=\chi E$
untuk memperoleh

$$\eqalignno{((\nu_1*\nu_2)*\nu_3)(E)
&=\iint\chi E(x+z)(\nu_1*\nu_2)(dx)\nu_3(dz)\cr
&=\iiint\chi E(x+y+z)\nu_1(dx)\nu_2(dy)\nu_3(dz)\cr
\noalign{\noindent (karena $x\mapsto \chi E(x+z)$ terukur Borel
untuk setiap $z$)}
&=\iint\chi E(x+y)\nu_1(dx)(\nu_2*\nu_3)(dy)\cr
\noalign{\noindent (karena $(x,y)\mapsto\chi E(x+y)$
terukur Borel, sehingga $y\mapsto\int\chi E(x+y)\nu_1(dx)$
terintegralkan terhadap $(\nu_2*\nu_3)$)}
&=(\nu_1*(\nu_2*\nu_3))(E).\cr}$$

\noindent Jadi $(\nu_1*\nu_2)*\nu_3$ dan $\nu_1*(\nu_2*\nu_3)$ bernilai sama pada
himpunan-himpunan Borel di $\BbbR^r$;  karena keduanya merupakan ukuran Radon,
keduanya harus identik.
}%end of proof of 257E

\leader{257F}{Teorema} Misalkan $\nu_1$ dan $\nu_2$ adalah ukuran Radon
berhingga total pada $\BbbR^r$ yang merupakan ukuran integral tak tentu
terhadap ukuran Lebesgue $\mu$.   Maka $\nu_1*\nu_2$ juga merupakan
ukuran integral tak tentu terhadap $\mu$;  jika $f_1$ dan $f_2$ masing-masing merupakan
turunan Radon--Nikod\'ym dari $\nu_1$, $\nu_2$, maka
$f_1*f_2$ merupakan turunan Radon--Nikod\'ym dari $\nu_1*\nu_2$.

\proof{ Berdasarkan 255H/255L, $f_1*f_2$ terintegralkan terhadap
$\mu$,
dengan $\int f_1*f_2d\mu=\nu_1(\BbbR^r)\nu_2(\BbbR^r)<\infty$, dan tentu saja
$f_1*f_2$ nonnegatif.   Jika
$E\subseteq\BbbR^r$ adalah suatu himpunan Borel,

$$\eqalignno{\int_Ef_1*f_2d\mu
&=\iint\chi E(x+y)f_1(x)f_2(y)\mu(dx)\mu(dy)\cr
\noalign{\noindent (255G)}
&=\iint\chi E(x+y)f_2(y)\nu_1(dx)\mu(dy)\cr
\noalign{\noindent (karena $x\mapsto\chi E(x+y)$ terukur Borel)}
&=\iint\chi E(x+y)\nu_1(dx)\nu_2(dy)\cr
\noalign{\noindent (karena $(x,y)\mapsto\chi E(x+y)$
terukur Borel, sehingga $y\mapsto\int\chi E(x+y)\nu_1(dx)$
terintegralkan terhadap $\nu_2$)}
&=(\nu_1*\nu_2)(E).\cr}$$

\noindent Jadi $f_1*f_2$ adalah turunan Radon--Nikod\'ym dari $\nu_1*\nu_2$ terhadap
$\mu$, berdasarkan 256J.
}%end of proof of 257F

\exercises{\leader{257X}{Latihan dasar $\pmb{>}$(a)}
%\spheader 257Xa
Misalkan $r\ge 1$ suatu bilangan bulat.   Misalkan $\delta_0$ adalah
ukuran Dirac pada $\BbbR^r$
yang terkonsentrasi di $0$.   Tunjukkan bahwa $\delta_0$ adalah ukuran probabilitas
Radon pada $\BbbR^r$ dan bahwa
$\delta_0*\nu=\nu$ untuk setiap $\nu$ yang merupakan ukuran Radon berhingga total
pada $\BbbR^r$.
%257A

\spheader 257Xb Misalkan $\mu$ dan $\nu$ adalah ukuran Radon berhingga total pada
$\BbbR^r$, dan $E$ sebarang himpunan yang terukur terhadap konvolusi keduanya $\mu*\nu$.
Tunjukkan bahwa $\int\mu(E-y)\nu(dy)$ terdefinisi
dalam $[0,\infty]$ dan sama dengan $(\mu*\nu)(E)$.
%257B

\spheader 257Xc Misalkan $\nu_1,\ldots,\nu_n$ adalah ukuran Radon berhingga total
pada $\BbbR^r$, dan misalkan $\nu$ adalah konvolusi
$\nu_1*\ldots*\nu_n$ (gunakan 257E untuk melihat bahwa ungkapan tanpa tanda kurung
seperti itu sah).   Tunjukkan bahwa

\Centerline{$\int h(x)\nu(dx)=\int\ldots\int
h(x_1+\ldots+x_n)\nu_1(dx_1)\ldots\nu_n(dx_n)$}

\noindent untuk setiap fungsi yang terintegralkan terhadap $\nu$, katakanlah $h$.
%257E

\spheader 257Xd Misalkan $\nu_1$ dan $\nu_2$ adalah ukuran Radon berhingga total
pada $\BbbR^r$, dengan dukungan $F_1$, $F_2$ (256Xf).   Tunjukkan bahwa
dukungan $\nu_1*\nu_2$ adalah $\overline{\{x+y:x\in F_1,\,y\in F_2\}}$.
%257E

\sqheader 257Xe Misalkan $\nu_1$ dan $\nu_2$ adalah ukuran Radon berhingga total
pada $\BbbR^r$, dan andaikan $\nu_1$ mempunyai turunan Radon--Nikod\'ym
$f$ terhadap
ukuran Lebesgue $\mu$.   Tunjukkan bahwa $\nu_1*\nu_2$ mempunyai
turunan Radon--Nikod\'ym $g$,
di mana $g(x)=\int f(x-y)\nu_2(dy)$ untuk $\mu$-hampir setiap $x\in\BbbR^r$.
%257F

\spheader 257Xf Misalkan $\nu_1$, $\nu_2$, $\nu_1'$ dan $\nu_2'$ adalah
ukuran Radon berhingga total pada $\BbbR^r$, dan andaikan bahwa $\nu_1'$, $\nu_2'$
masing-masing kontinu mutlak terhadap $\nu_1$, $\nu_2$.   Tunjukkan bahwa
$\nu_1'*\nu_2'$ kontinu mutlak terhadap $\nu_1*\nu_2$.
%257F

\leader{257Y}{Latihan lanjutan (a)}
Misalkan $M$ adalah ruang fungsional aditif
terhitung yang didefinisikan pada aljabar $\Cal B$ dari himpunan-himpunan bagian Borel
$\Bbb R$, dengan norma $\|\nu\|=|\nu|(\Bbb R)$ (lihat 231Yh).   (i) Tunjukkan
bahwa terdapat operator bilinear tunggal $*:M\times M\to M$ sedemikian sehingga
$(\mu_1\restr\Cal B)*(\mu_2\restr\Cal B)=(\mu_1*\mu_2)\restr\Cal B$ untuk
semua ukuran Radon berhingga total $\mu_1$, $\mu_2$ pada $\Bbb R$.  (ii)
Tunjukkan bahwa $*$ komutatif dan asosiatif.   (iii) Tunjukkan bahwa
$\|\nu_1*\nu_2\|\le\|\nu_1\|\|\nu_2\|$ untuk semua $\nu_1$, $\nu_2\in M$, sehingga
$M$ merupakan aljabar Banach di bawah perkalian ini.   (iv) Tunjukkan bahwa
$M$ mempunyai identitas perkalian.   (v) Tunjukkan bahwa $L^1(\mu)$ dapat
dipandang sebagai subaljabar tertutup dari $M$, dengan $\mu$ ukuran Lebesgue
pada $\Bbb R$ (bdk.\ 255Xi).
%257A

\spheader 257Yb Kita katakan bahwa suatu {\bf ukuran Radon
pada $\ocint{-\pi,\pi}$} adalah ukuran lengkap $\nu$ pada $\ocint{-\pi,\pi}$
sedemikian sehingga (i) setiap himpunan bagian Borel dari
$\ocint{-\pi,\pi}$ termasuk dalam domain $\Sigma$
dari $\nu$; (ii) untuk setiap $E\in\Sigma$
terdapat himpunan Borel $E_1$, $E_2$ sedemikian sehingga $E_1\subseteq E\subseteq E_2$
dan $\nu(E_2\setminus E_1)=0$; dan (iii) setiap himpunan bagian kompak
dari $\ocint{-\pi,\pi}$ mempunyai ukuran berhingga.   Tunjukkan bahwa untuk sebarang dua
ukuran Radon berhingga total $\nu_1$, $\nu_2$ pada $\ocint{-\pi,\pi}$
terdapat tepat satu ukuran Radon berhingga total $\nu$ pada
$\ocint{-\pi,\pi}$ sedemikian sehingga

\Centerline{$\int h(x)\nu(dx)=\int h(x+_{2\pi}y)\nu_1(dx)\nu_2(dy)$}

\noindent untuk setiap fungsi yang terintegralkan terhadap $\nu$, katakanlah $h$, dengan $+_{2\pi}$
didefinisikan seperti dalam 255Ma.
}%end of exercises

\endnotes{
\Notesheader{257}
Tentu saja konvolusi fungsi dan konvolusi ukuran berkaitan sangat
erat;  kaitan yang jelas ialah 257F, tetapi kesepadanan
antara 255G dan 257B juga sangat nyata.   Pada dasarnya, keduanya memberi kita
pengertian konvolusi $u*v$ yang sama ketika $u$, $v$ adalah anggota positif
dari $L^1$ dan $u*v$ ditafsirkan dalam $L^1$, bukan sebagai suatu fungsi
(257Ya).   Namun
kita harus menelaah jauh lebih dalam daripada argumen-argumen di sini untuk
menemukan gagasan dalam teori konvolusi ukuran yang bersesuaian dengan
hasil seperti 255K.   Saya akan kembali ke pertanyaan-pertanyaan jenis ini dalam
\S444 dalam Jilid 4.

Semua teorema dalam bagian ini dapat diperluas ke sebarang grup topologis
abelian Hausdorff yang kompak lokal;  tetapi untuk tingkat keumuman demikian
kita memerlukan gagasan yang jauh lebih maju (lihat \S444), dan untuk saat ini saya
hanya memberikan
saran dalam 257Yb agar Anda mencoba mengadaptasi gagasan-gagasan di sini ke
$\ocint{-\pi,\pi}$ atau $S^1$.
}%end of comment

\discrpage
 
